\section{CMake 说明}

本章逐段解释本工程的 CMake 配置，帮助理解「为什么这样写」。

\subsection{工具链文件 cmake/arm-none-eabi-gcc.cmake}

\begin{lstlisting}[style=generic]
set(CMAKE_SYSTEM_NAME Generic)                 # 目标系统：裸机，无操作系统
set(CMAKE_SYSTEM_PROCESSOR arm)
set(CMAKE_TRY_COMPILE_TARGET_TYPE STATIC_LIBRARY)  # 交叉编译时只能"试编译"，不能"试运行"

set(CMAKE_C_COMPILER     arm-none-eabi-gcc)
set(CMAKE_ASM_COMPILER   arm-none-eabi-gcc)
set(CMAKE_CXX_COMPILER   arm-none-eabi-g++)
set(CMAKE_OBJCOPY        arm-none-eabi-objcopy)
set(CMAKE_OBJDUMP        arm-none-eabi-objdump)
set(CMAKE_SIZE           arm-none-eabi-size)
\end{lstlisting}

要点：

\begin{itemize}
  \item \texttt{CMAKE\_SYSTEM\_NAME Generic}：告诉 CMake 这是裸机/无系统目标，
        不要链接宿主机（macOS/Linux）的标准库。
  \item \texttt{CMAKE\_TRY\_COMPILE\_TARGET\_TYPE STATIC\_LIBRARY}：
        交叉编译时，CMake 的编译器探测会尝试\textbf{运行}一个测试程序，
        但 arm 程序无法在宿主机上运行；设成静态库让探测只编译不运行，
        否则会报 \texttt{cannot run test program while cross compiling}。
  \item 编译器名不带路径：依赖 \texttt{PATH}。若想更鲁棒，可写成绝对路径。
\end{itemize}

\subsection{CMakeLists.txt 逐段解析}

\subsubsection{项目与语言}

\begin{lstlisting}[style=generic]
cmake_minimum_required(VERSION 3.24)
project(MAKE LANGUAGES C ASM)
set(CMAKE_C_STANDARD 99)
set(CMAKE_C_STANDARD_REQUIRED ON)
\end{lstlisting}

\texttt{LANGUAGES C ASM} 声明只用 C 和汇编（不启用 C++），
\texttt{ASM} 语言让 CMake 能直接编译 \texttt{.s} 启动文件。
\texttt{CMAKE\_C\_STANDARD} 指定 C 标准（本工程当前用 C99）。

\subsubsection{CPU 选项}

\begin{lstlisting}[style=generic]
set(CPU_FLAGS
    -mcpu=cortex-m4          # 目标内核
    -mthumb                  # 使用 Thumb-2 指令集（Cortex-M 只支持 Thumb）
    -mfpu=fpv4-sp-d16        # 单精度 FPU
    -mfloat-abi=hard         # 硬浮点 ABI（用 FPU 寄存器传参）
)
\end{lstlisting}

这些选项在\textbf{编译和链接}时都要加上（\texttt{target\_compile\_options} 与
\texttt{target\_link\_options} 都引用了 \texttt{CPU\_FLAGS}），
否则链接阶段会因浮点 ABI 不匹配而报错。

\subsubsection{目标与源文件}

\begin{lstlisting}[style=generic]
add_executable(${TARGET}
    ${CMAKE_SOURCE_DIR}/User/src/main.c
    ${CMAKE_SOURCE_DIR}/Core/src/system_stm32f4xx.c
    ${CMAKE_SOURCE_DIR}/startup/startup_stm32f40_41xxx.s
    ${CMAKE_SOURCE_DIR}/Driver/STM32F4xx_StdPeriph_Driver/src/misc.c
    ...
)
\end{lstlisting}

这里只显式列出了当前用到的几个外设库源文件（GPIO、RCC、SYSCFG、MISC）。
如果以后用到其它外设，把对应 \texttt{stm32f4xx\_*.c} 加进来即可，
\textbf{不必}把整个 \texttt{Driver/src/} 全部编译（省时间、省体积）。

\subsubsection{头文件路径与 SYSTEM}

\begin{lstlisting}[style=generic]
# 自己的头文件：正常警告
target_include_directories(${TARGET} PRIVATE
    ${CMAKE_SOURCE_DIR}/User/inc
    ${CMAKE_SOURCE_DIR}/Core/inc
)

# 厂商头文件：SYSTEM，忽略其内部告警
target_include_directories(${TARGET} SYSTEM PRIVATE
    ${CMAKE_SOURCE_DIR}/Driver/CMSIS/Include
    ${CMAKE_SOURCE_DIR}/Driver/CMSIS/Device/ST/STM32F4xx/Include
    ${CMAKE_SOURCE_DIR}/Driver/STM32F4xx_StdPeriph_Driver/inc
)
\end{lstlisting}

\texttt{SYSTEM} 关键字等价于给编译器传 \texttt{-isystem}，
使这些目录里的头文件被视为「系统头文件」，其内部告警被抑制。
这是处理第三方头文件告警的标准做法。

\subsubsection{编译宏定义}

\begin{lstlisting}[style=generic]
target_compile_definitions(${TARGET} PRIVATE
    STM32F40_41xxx          # 芯片型号宏，让 stm32f4xx.h 选择对应寄存器组
    USE_STDPERIPH_DRIVER    # 让 stm32f4xx.h 包含 stm32f4xx_conf.h
)
\end{lstlisting}

这两个宏是 ST 头文件的「开关」：不定义 \texttt{STM32F40\_41xxx} 会导致
\texttt{stm32f4xx.h} 报错 \texttt{Please select the target device}；
不定义 \texttt{USE\_STDPERIPH\_DRIVER} 则外设库头文件不会被包含。

\subsubsection{链接选项}

\begin{lstlisting}[style=generic]
target_link_options(${TARGET} PRIVATE
    ${CPU_FLAGS}
    -T${LD_SCRIPT}           # 指定链接脚本
    -nostdlib                # 不链接标准 C 库
    -nostartfiles            # 不用编译器自带的启动文件
    -Wl,--gc-sections        # 丢弃未使用的段（减体积）
    -Wl,-Map=${CMAKE_BINARY_DIR}/${TARGET}.map  # 生成 map 文件
)
\end{lstlisting}

\begin{itemize}
  \item \texttt{-nostdlib} / \texttt{-nostartfiles}：裸机关键选项，
        启动流程完全由 \texttt{startup/*.s} 自己负责。
  \item \texttt{--gc-sections}：配合编译选项 \texttt{-ffunction-sections
        -fdata-sections} 才能生效；本工程目前未显式加这两个编译选项，
        若想进一步瘦身，可在 \texttt{target\_compile\_options} 里补上。
\end{itemize}

\subsubsection{后处理：生成 .bin/.hex}

\begin{lstlisting}[style=generic]
add_custom_command(TARGET ${TARGET} POST_BUILD
    COMMAND ${CMAKE_OBJCOPY} -O binary $<TARGET_FILE:${TARGET}> ${BIN}
    COMMAND ${CMAKE_OBJCOPY} -O ihex   $<TARGET_FILE:${TARGET}> ${HEX}
    COMMAND ${CMAKE_SIZE}   $<TARGET_FILE:${TARGET}>
    ...
)
\end{lstlisting}

链接完成后自动调用 \texttt{objcopy} 生成 \texttt{.bin} / \texttt{.hex}，
并打印各段大小（text/data/bss）。

\subsection{CMakePresets.json}

本工程根目录下有一个 \texttt{CMakePresets.json}（目前为空）。预设的作用是
把「生成器 + 工具链 + 选项」打包成一个名字，之后一行命令即可配置：

\begin{lstlisting}[style=generic]
{
  "version": 3,
  "configurePresets": [
    {
      "name": "arm",
      "generator": "Ninja",
      "toolchainFile": "${sourceDir}/cmake/arm-none-eabi-gcc.cmake",
      "binaryDir": "${sourceDir}/build"
    }
  ]
}
\end{lstlisting}

配置时即可简化为 \texttt{cmake --preset arm}，构建用
\texttt{cmake --build --preset arm}。
